\documentclass[12pt]{article}

% Core math
\usepackage{amsmath, amssymb, amsthm}
\usepackage{stmaryrd} % \llbracket, \rrbracket

% Tables with paragraph-style cells
\usepackage{tabularx}

% Diagrams
\usepackage{tikz-cd}
\usepackage{tikz}

% References
\usepackage{hyperref}
\usepackage{cleveref}

% Graphics
\usepackage{graphicx}

% Code listings
\usepackage{listings}

% Page layout
\usepackage[margin=1in]{geometry}

% GrokRxiv sidebar (uses modern LaTeX shipout/background hook)
\usepackage{xcolor}

% Theorem environments
\newtheorem{theorem}{Theorem}[section]
\newtheorem{proposition}[theorem]{Proposition}
\newtheorem{lemma}[theorem]{Lemma}
\newtheorem{corollary}[theorem]{Corollary}
\theoremstyle{definition}
\newtheorem{definition}[theorem]{Definition}
\newtheorem{example}[theorem]{Example}
\theoremstyle{remark}
\newtheorem{remark}[theorem]{Remark}

% GrokRxiv sidebar definition (first page only)
\definecolor{grokgray}{RGB}{110,110,110}
\AddToHook{shipout/background}{%
  \ifnum\value{page}=1
    \begin{tikzpicture}[remember picture, overlay]
      \node[
        rotate=90,
        anchor=south,
        font=\Large\sffamily\bfseries\color{grokgray},
        inner sep=0pt
      ] at ([xshift=38pt, yshift=0.52\paperheight]current page.south west)
      {GrokRxiv:2026.05.naive-perspective-numbers\quad
       [\,math.HO\,]\quad
       03 May 2026};
    \end{tikzpicture}
  \fi
}

% Listings styling
\lstset{
  basicstyle=\ttfamily\small,
  columns=fullflexible,
  breaklines=true,
  showstringspaces=false,
  frame=single,
  framesep=4pt
}

\title{The Naive Perspective: \\ Numbers as Symbols and Quantities \\
\large Paper I of the Univalent Correspondence Series}

\author{YonedaAI Research \\
\textit{Univalent Correspondence Working Group} \\
\textit{Independent Research Collective} \\
\texttt{research@yonedaai.com}}

\date{3 May 2026}

\begin{document}
\maketitle

\begin{abstract}
This is the first paper in a six-part series, plus synthesis, on what we call
the \emph{Univalent Correspondence}: six perspectives on what a natural number
``is''. Across the series we proceed from the naive level (this paper),
through set-theoretic encoding, universal property, the Yoneda perspective,
Homotopy Type Theory (HoTT), and the categorical/structural view, before
unifying the latter four under Voevodsky's univalence axiom. The present
paper occupies the \emph{naive} or pre-formal level, where a number is at
once a symbol (``58'', \textsf{LVIII}, \texttt{||||\,||||}) and a quantity
(the cardinality of a counted collection). We treat numerals as syntax and
the underlying number as semantics, and we trace this distinction through
the historical arc that runs from Mill's empiricism, Husserl's phenomenology
of arithmetic, and Frege's critique of psychologism, into Benacerraf's
identification problem. We argue that the naive level is simultaneously
\emph{insufficient}---because of symbol/object confusion, the multiplicity
of equally adequate representations, and a foundational underdetermination
that already foreshadows the structuralist response---and \emph{indispensable},
because it is the interface where embodied counting, finger arithmetic,
and culturally transmitted numeral systems anchor mathematical content in
human cognition. We close by showing that the naive level is precisely
where the structuralist's question, ``what role does this object play?'',
first arises, thereby motivating each of the five formal levels developed
in subsequent papers. We also provide a small Haskell artifact demonstrating
the syntax-semantics distinction by parsing four numeral systems (Hindu-Arabic,
Roman, tally, Babylonian sexagesimal) into a single structural normal form.
\end{abstract}

\tableofcontents
\bigskip

\section{Introduction}
\label{sec:introduction}

\subsection{The question}

What is the number~58? On the page in front of you it is two glyphs, ``5''
and ``8''. In a Roman inscription it is the four-glyph sequence
\textsf{LVIII}. On a tally sheet kept by a shepherd it is fifty-eight scratches.
In Babylonian sexagesimal cuneiform it is a single wedge cluster meaning
``58''. Each of these is unambiguously a way of writing 58, and yet they share
no graphical features, no spelling, no length. Whatever 58 \emph{is}, it
cannot be any one of these inscriptions, for it survives translation between
all of them.

This paper is about the level of mathematical reflection at which one notices
this fact for the first time. We call this the \emph{naive} level, in the
technical, non-pejorative sense of being prior to formal axiomatization. It
is the level of the schoolchild learning to count, of the merchant tallying
goods, of the cuneiform scribe, and of every adult human who knows the
arithmetical facts of daily life without ever having read a definition of
the natural numbers. We will argue that this level is, at the same time,
deeply inadequate as a foundation and absolutely indispensable as a
phenomenology. The tension between these two claims sets up the entire
remainder of the series.

\subsection{Place in the series}

The series of which this paper is the first instalment is organized around
the following six-row table, which we will call the \emph{columns table}
throughout:

\begin{center}
\begin{tabularx}{\textwidth}{l|X}
\textbf{Level} & \textbf{What 58 \emph{is}} \\ \hline
Naive & A symbol; a quantity \\
Set-theoretic (von Neumann) & $\{0, 1, \ldots, 57\}$ \\
Universal property & The 58th iteration of $\mathrm{succ}$ in any initial successor structure \\
Yoneda & The representable functor $\mathrm{Hom}(-, 58)$ \\
HoTT & A canonical term in the inductive type $\mathbb{N}$, up to path equivalence \\
Categorical / structural & An invariant under structure-preserving morphisms of arithmetic \\
\end{tabularx}
\end{center}

Subsequent papers will treat rows two through six in turn, and the synthesis
paper will argue that rows three through six are, under the univalence
axiom, literally---not metaphorically---the same theory viewed through
different windows.

The job of \emph{this} paper is to establish row one with care, both as a
positive description of how mathematical content is normally encountered
and as a critical diagnosis of why we cannot stop there.

\subsection{Outline}

\Cref{sec:numerals-syntax} surveys numeral systems, treating them as syntax,
and articulates the syntax--semantics distinction in the form it takes at the
naive level. \Cref{sec:counting-quantity} develops the dual aspect---numbers
as quantities arising from counting and grouping---and connects this to
embodied cognition. \Cref{sec:pre-fregean} traces the pre-Fregean intellectual
history through Mill, Husserl, and Frege, with attention to Frege's critique
of psychologism and his Sinn/Bedeutung distinction. \Cref{sec:insufficient}
states three convergent reasons the naive level is insufficient as a
foundation: the symbol/object confusion, multiplicity of representations,
and Benacerraf's identification problem (anticipated already at the naive
level). \Cref{sec:indispensable} argues conversely that the naive level is
indispensable: counting, finger arithmetic, and embodied schemas are
constitutive of the meaning of arithmetic, not optional pedagogical scaffolding.
\Cref{sec:bridge} shows that the naive level is the natural locus for the
structuralist's question, ``what role does this object play?'', and that
this question is what every subsequent paper formalizes.
\Cref{sec:results} states the principal results of the paper as theorems---a
syntax/semantics adequacy lemma, a no-canonical-numeral theorem, and a
naive Benacerraf principle. \Cref{sec:haskell} describes a small Haskell
artifact that parses several numeral systems and reduces them all to a
single structural normal form. \Cref{sec:discussion,sec:conclusion} reflect
on connections and future work.

\section{Mathematical Framework}
\label{sec:framework}

This paper is unusual within the series in that its mathematical framework
is deliberately minimal: the entire point is to operate \emph{below} the
threshold of axiomatic formalization. Nevertheless, we will need just enough
formal language to say precisely what numerals are, what counting is, and
what it means for two numerals to denote the same number. We collect this
machinery here.

\subsection{Numerals as words over a finite alphabet}

\begin{definition}[Numeral system]
\label{def:numeral-system}
A \emph{numeral system} is a triple $\mathcal{N} = (\Sigma, L, \llbracket\,
\cdot\, \rrbracket)$ where:
\begin{enumerate}
\item $\Sigma$ is a finite alphabet of \emph{glyphs};
\item $L \subseteq \Sigma^*$ is a (possibly partial) language of well-formed
\emph{numerals}; and
\item $\llbracket\, \cdot\, \rrbracket : L \to \mathbb{N}$ is a
\emph{denotation function} mapping each well-formed numeral to a natural
number.
\end{enumerate}
We do not, at this stage, presuppose any particular formalization of
$\mathbb{N}$; the codomain is the pretheoretic collection of natural
numbers as understood by the practitioners of the system.
\end{definition}

\begin{example}[Hindu-Arabic decimal positional system]
\label{ex:hindu-arabic}
$\Sigma_{10} = \{0, 1, 2, \ldots, 9\}$, $L_{10}$ is the set of finite
non-empty strings over $\Sigma_{10}$ with no leading $0$ except for the
string ``$0$'' itself, and
$\llbracket d_{n-1} d_{n-2} \cdots d_1 d_0 \rrbracket = \sum_{i=0}^{n-1}
d_i \cdot 10^i$.
\end{example}

\begin{example}[Tally]
\label{ex:tally}
$\Sigma_{|} = \{\,|\,\}$, $L_{|} = \{\,|\,\}^*$ (all finite strings),
$\llbracket s \rrbracket = $ the length of $s$. We will sometimes use
``gate-five'' tally with a closing diagonal stroke; this changes
$\Sigma_{|}$ but not the denotation.
\end{example}

\begin{example}[Roman]
\label{ex:roman}
$\Sigma_R = \{\,\textsf{I}, \textsf{V}, \textsf{X}, \textsf{L},
\textsf{C}, \textsf{D}, \textsf{M}\,\}$ with the usual subtractive
constraints determining $L_R$, and the standard semantics yielding
$\llbracket \textsf{LVIII}\rrbracket = 58$.
\end{example}

\begin{example}[Babylonian sexagesimal]
\label{ex:babylonian}
A base-60 positional system using two cuneiform glyphs (one for ones, one
for tens) to compose digits 1--59, with positional weights $60^i$. The
historical absence of a true zero glyph in early Babylonian usage means
$L$ is partial and ambiguity is resolved by context; modern transliterations
add an explicit zero placeholder.
\end{example}

\subsection{Equivalence of numerals}

\begin{definition}[Numeral equivalence]
\label{def:numeral-equivalence}
Two numerals $w_1 \in L_1$ and $w_2 \in L_2$ in possibly different numeral
systems $\mathcal{N}_1, \mathcal{N}_2$ are \emph{equivalent}, written
$w_1 \sim w_2$, iff $\llbracket w_1 \rrbracket_1 = \llbracket w_2
\rrbracket_2$.
\end{definition}

\begin{remark}
\label{rem:symbol-vs-number}
Already in \cref{def:numeral-equivalence}, we have implicitly committed to
a substantive thesis: that there is some target codomain in which
$\llbracket w_1 \rrbracket_1$ and $\llbracket w_2 \rrbracket_2$ can be
\emph{compared}. A consistently naive theorist who refuses any abstraction
beyond the inscriptions themselves cannot even state when two numerals
``mean the same thing''. The bare existence of \cref{def:numeral-equivalence}
is therefore evidence that the naive level is not internally stable: as
soon as one wishes to compare across numeral systems, one is reaching
beyond pure syntax.
\end{remark}

\subsection{Counting as a function}

\begin{definition}[Counting function]
\label{def:counting}
Let $X$ be a finite set whose elements are physically or conceptually
individuable. A \emph{counting} of $X$ is a bijection $c : \{1, 2, \ldots,
n\} \to X$ for some $n \in \mathbb{N}$. The number $n$ is called the
\emph{cardinality} of $X$.
\end{definition}

\begin{remark}
The \emph{availability} of \cref{def:counting} again presupposes more
than the strictly naive theorist allows: it presupposes the natural numbers
as an indexing device. We will return to this regress in
\cref{sec:insufficient}.
\end{remark}

\section{Numerals as Syntax}
\label{sec:numerals-syntax}

\subsection{Positional, additive, and tally systems}

The world's historical numeral systems can be roughly classified along two
orthogonal axes: \emph{base} (the size of the radix used to bundle
quantities) and \emph{principle} (positional vs.\ additive vs.\ pure tally).

\paragraph{Tally.} The pure tally is base-1 in the sense that each unit
contributes additively and there is no bundling. It is the simplest
syntactic system one can imagine. Bone artifacts such as the Lebombo bone
(c.\ 42{,}000 BCE) and the Ishango bone (c.\ 18{,}000--20{,}000 BCE) are
typically interpreted as tallies.

\paragraph{Additive systems.} Roman, Egyptian, and Greek alphabetic
numerals are additive: distinct glyphs encode distinct values, and the
value of a numeral is obtained (with local subtractive corrections) by
summing the values of its glyphs. Such systems compress the tally---one
glyph for ten, another for fifty, another for a hundred---without yet
exploiting position.

\paragraph{Positional systems.} The Babylonian sexagesimal system,
Mesoamerican vigesimal systems, and the Hindu-Arabic decimal system are
positional: a small alphabet of glyphs is reused at each ``place,'' and
the place itself contributes a power of the base. Positional systems are
exponentially more compact than additive ones in the length of the
numeral required to denote a given number.

\begin{proposition}[Compression by positionality]
\label{prop:positional-compression}
Let $b \geq 2$ be a base. A positional numeral system of base $b$ requires
a numeral of length $\Theta(\log_b n)$ to denote $n$, whereas a pure tally
system requires length $n$, and an additive system with $k$ glyphs of
distinct values requires length $\Theta(n / v_{\max})$ where $v_{\max}$
is the largest glyph value (asymptotically linear in $n$).
\end{proposition}

\begin{proof}[Proof sketch]
Standard. For positional, $n$ has $\lfloor \log_b n \rfloor + 1$ digits.
For tally, length equals $n$. For additive, the optimal greedy
representation uses $\Theta(n / v_{\max})$ tokens of the largest available
glyph, plus $O(1)$ corrections.
\end{proof}

\subsection{Syntax versus semantics, naively stated}

\begin{definition}[Naive syntax/semantics distinction]
\label{def:naive-ss}
Fix a numeral system $\mathcal{N} = (\Sigma, L, \llbracket\,\cdot\,\rrbracket)$.
Then:
\begin{itemize}
\item the \emph{syntactic} content of a numeral $w$ is $w$ itself, the
finite string of glyphs;
\item the \emph{semantic} content is the value $\llbracket w \rrbracket$.
\end{itemize}
The fundamental observation of the naive level is that distinct syntactic
strings can have identical semantic content---even \emph{within} one numeral
system (e.g.\ \textsf{IIII} vs.\ \textsf{IV} in pre-medieval vs.\ medieval
Roman usage)---and \emph{a fortiori} across systems.
\end{definition}

We emphasize that \cref{def:naive-ss} is a \emph{weak} version of the
syntax/semantics distinction. The semantic codomain is the natural numbers
treated as a primitive concept; we have not (yet) said what natural numbers
are. The remaining papers in the series can each be read as offering a
candidate filling of this gap.

\subsection{Multiplicity and underdetermination}

The fact that there are infinitely many numeral systems for the same
underlying $\mathbb{N}$ is not a defect of any one system; it is a feature
of the situation. But it generates a diagnostic question that returns
again and again throughout the series: \emph{which features of the
representation are essential to the number, and which are accidents of
the inscription?} Phrasing this question already commits one to thinking
of numbers as something \emph{other than} their inscriptions. We will see
in \cref{sec:insufficient} that this is the first crack in the naive
edifice.

\section{Counting and Quantity}
\label{sec:counting-quantity}

\subsection{The dual aspect: number as quantity}

The other face of the naive level is the use of numbers to express
\emph{quantity}. A flock has 58 sheep; a sentence has 58 letters; a year
has 365 days. The number-as-quantity face is older than any numeral
system: archeological and ethnographic evidence of one-to-one correspondence
between counted items and tally marks predates positional notation by tens
of thousands of years.

\begin{definition}[Naive cardinality, repeated for emphasis]
\label{def:naive-cardinality}
The \emph{cardinality} of a finite collection $X$ is the result of placing
the elements of $X$ in one-to-one correspondence with an initial segment
of the count sequence. In particular, two collections have the same
cardinality iff they can be put into one-to-one correspondence with each
other. We do not, at this stage, presuppose anywhere that this number
``inhabits.''
\end{definition}

\begin{example}[Hume's principle, anticipated]
\label{ex:hume}
The following is sometimes called Hume's principle (after a remark in
Hume's \emph{Treatise} I.iii.1): the number of $F$s equals the number of
$G$s iff the $F$s and the $G$s can be put in one-to-one correspondence.
This is the naive level's characterization of cardinal identity. Frege
famously elevated it to a foundational role in his \emph{Grundlagen}
(1884), and neologicism (Wright, Hale, Heck) takes it as definitionally
sufficient for arithmetic.
\end{example}

\subsection{The count sequence and embodied cognition}

The cognitive science literature documents that human and pre-linguistic
infant counting relies on at least two distinct systems: a precise
\emph{small number system} (subitizing, accurate up to about 4) and an
approximate large number system whose Weber-fraction discriminability is
roughly logarithmic. Numerals expand the precise system arbitrarily by
piggy-backing on language and on culturally transmitted iteration schemes
(``next, next, next''). Lakoff and N\'u\~nez argue that the count sequence
itself is a conceptual metaphor grounded in embodied schemas: motion along
a path, the iteration of grasping motions, the sequential pointing of a
finger. We do not endorse all the metaphysics in their account, but we
follow them on one thesis that we will need: \emph{the count sequence is
not optional}. There is no human acquisition of arithmetic that bypasses
the embodied count.

\begin{remark}[Finger arithmetic]
\label{rem:fingers}
A working adult who has internalized the multiplication table to 9$\times$9
typically retains traces of a finger-based representation, accessible as
imagery if not in overt behavior. The Chinese ``Chisanbop'' technique and
the Vedic Indian finger-counting traditions are explicit, mature systems
in which finger configurations \emph{are} the numerals. They illustrate
that the line between syntax (numerals as inscriptions) and embodied
quantity (numbers as counted) is not where one might naively think it is:
fingers are inscriptions on the body.
\end{remark}

\subsection{Counting as proto-induction}

\begin{proposition}[Naive induction principle]
\label{prop:naive-induction}
Any property $P$ that (i) holds of an initial segment of the count
sequence and (ii) is preserved when one passes from the current count to
the next, holds of every count one will reach.
\end{proposition}

\begin{proof}[Discussion]
At the naive level, this is not a theorem; it is the implicit principle
that governs counting. ``Once you've counted that high, you can keep
going.'' Formalizing this principle is the work of the universal-property
paper (Paper III) and the HoTT paper (Paper V). What we record here is
that the principle is already \emph{operative} at the naive level, in the
sense that competent counters use it without knowing they are using it.
\end{proof}

\subsection{Quantity does not entail abstract object}

A subtle but important observation. The naive use of numbers as quantities
does not, by itself, commit one to numbers being objects. ``There are 58
sheep'' may, on a deflationary reading, mean only that each of the 58
sheep is a sheep and no further sheep is. This deflationary reading is
internally consistent and, we believe, captures something real about how
numbers function in the most everyday cases. But it cannot survive the
move to \emph{arithmetic}---to facts of the form $58 = 2 \cdot 29$---without
some non-trivial reification. We will return to this point in
\cref{sec:insufficient}.

\section{The Pre-Fregean and Early Fregean Background}
\label{sec:pre-fregean}

\subsection{Mill's empiricism}

John Stuart Mill, in his \emph{System of Logic} (1843), Book II, Chapter
VI, defended the view that arithmetical propositions are inductive
generalizations from physical experience. ``$2 + 1 = 3$'' is, on his view,
a generalization observed to hold whenever one combines a pair of pebbles
with a single pebble. This view has the merit of explaining the
applicability of arithmetic to the world, but it suffers from at least
two well-known defects:
\begin{enumerate}
\item Arithmetical propositions appear necessary; inductive generalizations
do not. We do not feel that ``$2+1=3$'' could fail in the way that ``every
swan is white'' could fail.
\item Arithmetical propositions apply to abstract and even imagined
collections in exactly the way they apply to physical collections. The
inductive base, on Mill's view, ought to be physical collections; the
extension to abstracta then becomes mysterious.
\end{enumerate}
Frege made the criticism of Mill central in his \emph{Grundlagen} (1884),
\S\S 7--10. We side with the orthodox post-Fregean reading on this, while
recording that contemporary cognitive science (small-number system,
approximate number system, neural representations of magnitude) gives
Mill's account a partial vindication \emph{at the level of psychology}:
something \emph{like} an empirical magnitude-sense is part of the
acquisition story, even if it is not part of the justification story.

\subsection{Husserl's phenomenology of arithmetic}

Edmund Husserl's \emph{Philosophie der Arithmetik} (1891) attempted to
ground arithmetic in mental acts of \emph{collective combination}: the
mental act of bringing items together as a multiplicity. Frege's
\emph{Review of Husserl} (1894) is one of the founding documents of
analytic philosophy, and its central charge is \emph{psychologism}: the
confusion of an objective content (a number, an arithmetical truth) with
the subjective act by which a thinker entertains it. Husserl himself
later moved away from this view; the mature Husserl (\emph{Logical
Investigations}, 1900--01) accepts a sharp ideal/real distinction.

\subsection{Frege's critique of psychologism}

Frege's critique can be summarized as a slogan: \emph{the truth that $7$
is prime does not depend on anyone's thinking it.} If it did, then by
varying the thinker one would vary the truth, which is absurd. The same
slogan re-purposed against the strict naive level says: \emph{the number
58 cannot be the inscription ``58'',} for one and the same number is
inscribed by many different glyph sequences.

\begin{remark}[Sinn and Bedeutung]
\label{rem:sinn-bedeutung}
Frege's later distinction (\emph{\"Uber Sinn und Bedeutung}, 1892)
between \emph{sense} (Sinn) and \emph{reference} (Bedeutung) is the
mature articulation of the syntax/semantics distinction we have been
developing. ``58'' and ``$2 \cdot 29$'' are different senses of the same
referent. The naive level treats sense and reference as identical (or at
least as undistinguished). Distinguishing them is one of the
preconditions of all the formal levels to come.
\end{remark}

\subsection{Frege's logicism, briefly}

Frege's positive program in the \emph{Grundlagen} and \emph{Grundgesetze}
was to define numbers as extensions of certain second-order concepts:
the number of $F$s is the extension of the concept ``equinumerous with
$F$''. Russell's paradox showed that the unrestricted comprehension
needed for such extensions is inconsistent. This is not the place to
recount the recovery program (Russell's type theory, Zermelo's set
theory, neologicism); we record only that Frege's specific reduction
fails, while his \emph{methodological} thesis---that arithmetic is
prior to and independent of psychology---is the orthodox view we take
forward into the rest of the series.

\section{Why the Naive Level is Insufficient}
\label{sec:insufficient}

We give three convergent arguments. They overlap, but each isolates a
distinct failure mode.

\subsection{The symbol/object confusion}

\begin{theorem}[Symbol/object confusion theorem]
\label{thm:symbol-object}
There is no consistent way, working purely at the naive level, to identify
a number with its inscription. In particular, if one assumes (i) that the
glyph sequence ``58'' \emph{is} the number 58, and (ii) that
$58 = \mathsf{LVIII}$, then one is forced to identify the glyph sequences
``58'' and ``\textsf{LVIII}'', which are manifestly distinct.
\end{theorem}

\begin{proof}
By transitivity of identity. If ``58'' $= 58$ and $58 = \mathsf{LVIII}$,
then ``58'' $= \mathsf{LVIII}$. But the strings ``58'' and
``\textsf{LVIII}'' are distinct strings of distinct glyphs over different
alphabets. Contradiction.
\end{proof}

The argument is elementary, but it has weight. It shows that the very
\emph{intelligibility} of equations like $58 = \mathsf{LVIII}$ requires
distinguishing the number from any one of its inscriptions. The naive
level is the level at which this distinction is not yet drawn; the rest
of the series is the work of drawing it carefully.

\subsection{Multiplicity of representations}

\begin{theorem}[No-canonical-numeral theorem]
\label{thm:no-canonical-numeral}
There is no privileged numeral system. Specifically, for any finite set
of numeral systems $\{\mathcal{N}_1, \ldots, \mathcal{N}_k\}$ each adequate
for representing every natural number, there is a further numeral system
$\mathcal{N}_{k+1}$, not equivalent to any $\mathcal{N}_i$ as a
syntactic object, that is also adequate.
\end{theorem}

\begin{proof}
Given $\mathcal{N}_1, \ldots, \mathcal{N}_k$, define $\mathcal{N}_{k+1}$
as the positional system over base $b' = \max_i |\Sigma_i| + 1$. By
construction, $|\Sigma_{k+1}| = b'$ exceeds the alphabet size of every
$\mathcal{N}_i$, so $\mathcal{N}_{k+1}$ is not equivalent as a syntactic
object to any of them. Adequacy is immediate from the standard properties
of positional notation.
\end{proof}

\Cref{thm:no-canonical-numeral} is a syntactic version of a more famous
theorem we will meet in Paper II:

\subsection{Naive Benacerraf}

\begin{theorem}[Naive Benacerraf principle]
\label{thm:naive-benacerraf}
Every numeral system $\mathcal{N}_i$ adequate for $\mathbb{N}$ is
``equally good'' from the standpoint of arithmetic. Hence no
$\mathcal{N}_i$ can be \emph{the} number system; numbers must be the
common semantic content shared by all adequate $\mathcal{N}_i$.
\end{theorem}

\begin{proof}[Proof sketch]
Let $\mathcal{N}_i, \mathcal{N}_j$ be adequate. For any
arithmetical operation $\circ$ definable in both, one verifies that
$\llbracket w_1 \circ w_2 \rrbracket_i = \llbracket w_1' \circ w_2'
\rrbracket_j$ whenever $w_1 \sim w_1'$ and $w_2 \sim w_2'$, where $\sim$
is the cross-system equivalence of \cref{def:numeral-equivalence}. This
witnesses that $\mathcal{N}_i$ and $\mathcal{N}_j$ are interconvertible
without arithmetical loss.
\end{proof}

\begin{remark}
This is essentially Benacerraf's 1965 \emph{identification problem}
``one floor down''. Benacerraf's problem asks why we should regard the
von Neumann encoding as canonical when the Zermelo encoding is equally
adequate. The naive Benacerraf principle observes the same multiplicity
already at the level of numeral systems, before set theory is even
brought in. Paper II will pursue Benacerraf at the set-theoretic level
proper.
\end{remark}

\subsection{Convergence: numbers cannot be naive things}

The three arguments above (symbol/object confusion, multiplicity of
representations, naive Benacerraf) jointly establish:

\begin{theorem}[Inadequacy of the naive level]
\label{thm:naive-inadequate}
A theory of arithmetic that identifies numbers with their inscriptions,
or equivalently with any single privileged numeral system, cannot account
for the existence of cross-system equations like $58 = \mathsf{LVIII}$.
Therefore, numbers are not naively given things; they are something
shared across all adequate inscriptions.
\end{theorem}

\begin{proof}
Combine \cref{thm:symbol-object,thm:no-canonical-numeral,thm:naive-benacerraf}.
\end{proof}

\Cref{thm:naive-inadequate} is the negative result of the paper: the
naive level cannot be the foundational level. The remainder of the
series is the constructive search for what numbers \emph{are}, given
that they are not their inscriptions.

\section{Why the Naive Level is Indispensable}
\label{sec:indispensable}

The negative results of the previous section are sometimes overread, in
the foundationalist literature, as showing that the naive level should be
\emph{discarded} once the formal levels are in place. We disagree, and
this section explains why.

\subsection{Indispensability arguments}

\begin{proposition}[Acquisition indispensability]
\label{prop:acquisition}
No human learns arithmetic without learning to count. The count sequence
and the elementary numerals are the entry point into the entire
mathematical edifice; no formal system of arithmetic has ever been
acquired by anyone without first traversing the naive level.
\end{proposition}

\begin{proof}[Empirical sketch]
This is an empirical claim, supported by the developmental psychology
literature (Piaget, Gelman \& Gallistel, Carey). Mathematics education
universally proceeds by way of numerals and counting; no curriculum
introduces the Peano axioms or the natural number object first.
\end{proof}

\begin{proposition}[Application indispensability]
\label{prop:application}
Every application of arithmetic to the empirical world proceeds via
counting. The bridge from the formal natural numbers to actual sheep,
electrons, or seconds is provided by an act of counting (or its
extension, measurement), which is a naive operation.
\end{proposition}

\begin{proof}[Discussion]
A formalization of arithmetic gives the structural facts (e.g.\ that
$58 = 2 \cdot 29$). It does not, by itself, tell us that there are 58
sheep in the field. The latter requires an act of counting that connects
the formal token to the empirical scene. Counting is irreducibly naive
in our sense.
\end{proof}

\begin{proposition}[Pedagogical indispensability]
\label{prop:pedagogical}
The introduction of any of the formal levels (universal property,
Yoneda, HoTT) presupposes facility with the naive level. The student
who is to be taught the universal property of $\mathbb{N}$ must already
know what $\mathbb{N}$ \emph{is intuitively}; the universal property
purports to characterize what is being introduced, and the
characterization is empty if the introduced object is not already
familiar.
\end{proposition}

\subsection{The naive level as semantic anchor}

The cumulative force of \cref{prop:acquisition,prop:application,%
prop:pedagogical} is that the naive level is what \emph{makes the formal
levels semantically non-empty}. A formal system is, considered abstractly,
a manipulation of inscriptions according to rules. What gives it
mathematical content is the connection to the naive practice it is
designed to capture. This connection is not optional; it is constitutive.

\begin{remark}[Compare: physical theories]
\label{rem:physics-compare}
This is similar in form to the role played in physics by operational
definitions of measurement. A formal theory of mechanics is empty
without the bridge laws that say what counts as a measurement of length,
of time, of mass. The naive level of arithmetic plays the analogous
bridging role between inscription-manipulation and quantitative content.
\end{remark}

\subsection{Insufficient and indispensable: not a paradox}

The two main claims of the paper---that the naive level is insufficient
(\cref{sec:insufficient}) and that it is indispensable
(\cref{sec:indispensable})---are not in tension. They occupy different
levels of analysis. The naive level is insufficient as a \emph{foundation}
because it cannot answer foundational questions; it is indispensable as
a \emph{phenomenology} because all foundational work begins from it and
returns to it. The whole series should be read as taking this dual stance
seriously: the formal levels do real foundational work, but they do so
by way of, not by replacement of, the naive level.

\section{Bridge to the Formal Levels: The Structuralist's Question}
\label{sec:bridge}

The diagnostic question that emerges from \cref{thm:naive-inadequate} is:
if a number is not any of its inscriptions, what \emph{is} it? At the
naive level, the natural answer is something like ``the number is what
all these inscriptions have in common.'' Pressed for what they have in
common, the naive theorist can offer only: the same place in the count
sequence; the same arithmetical role.

\begin{definition}[The structuralist's question]
\label{def:structuralist-question}
The \emph{structuralist's question}, asked of a putative mathematical
object $X$, is: \emph{What role does $X$ play in the structure that
contains it?}
\end{definition}

\begin{proposition}[The structuralist's question first arises naively]
\label{prop:structuralist-naive}
The structuralist's question is the natural follow-up to
\cref{thm:naive-inadequate} as encountered at the naive level. Numbers
are what their inscriptions share; what their inscriptions share is the
arithmetical role; therefore, numbers are arithmetical roles.
\end{proposition}

This formulation is, of course, deeply unsatisfying as it stands. It
trades on intuitive notions of ``role'' and ``share'' that have not been
formalized. The five subsequent papers in the series can be read as
five different formalizations of \cref{prop:structuralist-naive}:

\begin{itemize}
\item \textbf{Paper II (Set-theoretic).} A number's role is encoded by
its membership profile in a chosen set-theoretic universe. Multiple
encodings (von Neumann, Zermelo) instantiate the same role; their
underdetermination is Benacerraf's identification problem.
\item \textbf{Paper III (Universal property).} A number's role is its
position in the initial successor structure $(\mathbb{N}, 0,
\mathrm{succ})$, characterized by its universal property. The
underdetermination of Paper II is dissolved.
\item \textbf{Paper IV (Yoneda).} A number's role is its representable
functor $\mathrm{Hom}(-, n)$. ``What role does $n$ play?'' becomes
``what natural transformations land at $n$?''.
\item \textbf{Paper V (HoTT).} A number's role is its identity (path)
type within the inductive type $\mathbb{N}$. Univalence makes the
role-identity equivalence native rather than merely modeled.
\item \textbf{Paper VI (Categorical / structural).} A number's role
is what is invariant under all structure-preserving morphisms of
arithmetic.
\end{itemize}

The synthesis paper argues that, under univalence, papers III through
VI are the same theory.

\begin{remark}[The naive level is recovered, not erased]
\label{rem:recovery}
A subtle point. The structuralist's question, in any of its formal
guises, gives a precise answer to ``what \emph{is} 58?'' But the answer
is licensed by the naive level: it is licensed because the structural
characterization \emph{recovers} the inscriptions and the counting practice
in the right way. A formalization that gave the right structural answer
but failed to recover the everyday use of ``58'' would be deficient as a
theory of \emph{arithmetic} (whatever else it might be). The naive level
is therefore not only a starting point but a constraint of adequacy on
all the formalizations.
\end{remark}

\section{Results}
\label{sec:results}

We collect the principal results of the paper for ease of reference.

\begin{theorem}[Symbol/object confusion]
The strict identification of a number with any one of its inscriptions
is inconsistent with the existence of cross-system equations.
(\cref{thm:symbol-object})
\end{theorem}

\begin{theorem}[No-canonical-numeral]
There is no privileged numeral system; for any finite set of adequate
numeral systems there is a further adequate system not equivalent as a
syntactic object to any of them. (\cref{thm:no-canonical-numeral})
\end{theorem}

\begin{theorem}[Naive Benacerraf principle]
All adequate numeral systems are interconvertible without arithmetical
loss; numbers are the shared semantic content of all of them.
(\cref{thm:naive-benacerraf})
\end{theorem}

\begin{theorem}[Inadequacy of the naive level]
A theory that identifies numbers with their inscriptions, or with any
single privileged numeral system, cannot underwrite cross-system
arithmetic. (\cref{thm:naive-inadequate})
\end{theorem}

\begin{proposition}[Indispensability of the naive level]
The naive level is constitutive of (a) acquisition of arithmetic, (b)
application of arithmetic, and (c) pedagogical introduction of formal
arithmetic. (Combination of \cref{prop:acquisition,prop:application,%
prop:pedagogical}.)
\end{proposition}

\begin{proposition}[Naive emergence of structuralism]
The structuralist's question, ``what role does $n$ play?'', arises
naturally from the inadequacy of the naive level and is the unifying
question of papers II through VI. (\cref{prop:structuralist-naive})
\end{proposition}

\section{The Haskell Artifact}
\label{sec:haskell}

To illustrate the syntax/semantics distinction concretely we provide a
small Haskell program that parses numerals in four systems---decimal,
Roman, tally, and a simplified Babylonian sexagesimal---and reduces them
all to a single structural normal form, namely a Peano natural number.

The program is intentionally minimal. Its purpose is not industrial
parsing but illustration of the following correspondence: distinct
syntactic constructors (the parsers) feed a single semantic constructor
(the normal form). What we model in code is exactly the diagram

\[
\begin{array}{ccc}
\text{numeral} & \xrightarrow{\;\text{parse}\;} & \text{semantic value (Peano natural)}
\end{array}
\]
with several different parsers and one shared codomain. We then test that
\textsf{LVIII}, \texttt{58}, a tally of $58$ strokes (i.e.\
\texttt{||||\dots||||} with $58$ vertical bars), and the sexagesimal
$(58, 0)$ all reduce to the same Peano value.

\subsection{Module structure}

The artifact is split into the following modules:
\begin{itemize}
\item \texttt{Naive.Peano}: the structural normal form (Peano naturals)
and basic arithmetic.
\item \texttt{Naive.Decimal}: a Hindu-Arabic decimal parser.
\item \texttt{Naive.Roman}: a Roman numeral parser handling subtractive
notation.
\item \texttt{Naive.Tally}: a tally parser.
\item \texttt{Naive.Babylonian}: a simplified base-60 parser.
\item \texttt{Main}: tests that all four parsers agree on a fixed sample.
\end{itemize}
The full source is at \texttt{src/01-naive/} in the project repository.

\subsection{The shared codomain}

\begin{lstlisting}[language=Haskell]
data Peano = Z | S Peano deriving (Eq, Show)

fromInt :: Int -> Peano
fromInt n
  | n <= 0    = Z
  | otherwise = S (fromInt (n - 1))

toInt :: Peano -> Int
toInt Z     = 0
toInt (S k) = 1 + toInt k

addP :: Peano -> Peano -> Peano
addP Z     y = y
addP (S x) y = S (addP x y)
\end{lstlisting}

This is the structural normal form into which all four parsers reduce.
We will return to this datum in Paper III, where we will recognize
\texttt{Peano} as an instance of the universal property of $\mathbb{N}$,
and in Paper V, where it appears as the inductive type defining
$\mathbb{N}$ in HoTT.

\subsection{Decimal, Roman, tally, Babylonian}

A sketch of each parser:

\begin{lstlisting}[language=Haskell]
parseDecimal :: String -> Maybe Peano
parseDecimal = fmap fromInt . readMaybe

romanValue :: Char -> Maybe Int
romanValue c = lookup c
  [('I',1),('V',5),('X',10),('L',50),('C',100),('D',500),('M',1000)]

parseRoman :: String -> Maybe Peano
parseRoman = fmap fromInt . roman . map romanValue
  where
    roman vs = case sequence vs of
      Nothing -> Nothing
      Just xs -> Just (foldr step 0 (zip xs (drop 1 xs ++ [0])))
    step (x,y) acc = if x < y then acc - x else acc + x

-- Tally: count only '|'; '/' is a grouping symbol that
-- contributes 0. Any other character fails.
parseTally :: String -> Maybe Peano
parseTally s
  | all (`elem` "|/") s = Just (fromInt (length (filter (=='|') s)))
  | otherwise           = Nothing

-- Babylonian: not a string parser. Given a pre-tokenized
-- list of (digit, place) pairs, return the denotation.
fromBabylonian :: [(Int,Int)] -> Maybe Peano
fromBabylonian places
  | all valid places = Just (fromInt (sum
      [d * 60 ^ p | (d,p) <- places]))
  | otherwise        = Nothing
  where valid (d, p) = 0 <= d && d < 60 && p >= 0
\end{lstlisting}

\subsection{The agreement test}

\begin{lstlisting}[language=Haskell]
main :: IO ()
main = do
  let dec   = parseDecimal   "58"
      rom   = parseRoman     "LVIII"
      tally = parseTally     (replicate 58 '|')
      bab   = fromBabylonian [(58, 0)]
  putStrLn $ "decimal:    " ++ show dec
  putStrLn $ "roman:      " ++ show rom
  putStrLn $ "tally:      " ++ show tally
  putStrLn $ "babylonian: " ++ show bab
  putStrLn $ "all equal?  " ++ show (and
    [ dec == rom, rom == tally, tally == bab ])
\end{lstlisting}

The point of the test is exactly the point of the paper: four glyphically
unrelated inscriptions reduce to the same structural object. The
Haskell type system enforces this: all four reduction functions have
codomain \texttt{Maybe Peano}, so equality of their outputs is
meaningful.

\begin{remark}[Notes on the artifact]
\label{rem:artifact-notes}
Two terminological points are worth flagging. First, \texttt{parseTally}
counts only \texttt{'|'} marks; the diagonal \texttt{'/'} is treated as
a grouping symbol contributing zero to the count. Hence, on this
parser, \texttt{"//|||"} and \texttt{"|||"} both denote 3. The choice
matches the paper's emphasis on denotation rather than well-formedness
of inscription. Second, \texttt{fromBabylonian} is not a string parser
but a denotation function on pre-tokenized \texttt{[(digit, place)]}
lists. We retain \texttt{parseBabylonian} as a legacy alias in the
shipped module. The full source treats this naming explicitly.
\end{remark}

\section{Discussion}
\label{sec:discussion}

\subsection{Connections to the rest of the series}

The naive level establishes the desiderata that each formal level must
meet. We summarize the form these desiderata take in each subsequent
paper:

\begin{itemize}
\item \emph{Paper II (set-theoretic)} must explain how the multiplicity
of inscriptions can be reduced to a single set-theoretic encoding---and
must confront the residual underdetermination at the next level (von
Neumann vs.\ Zermelo encoding).
\item \emph{Paper III (universal property)} must show that the
structural role identified at the naive level is captured by initiality
in the category of successor structures; in particular, the unique
recursion morphism must reproduce the everyday operations the naive
level uses.
\item \emph{Paper IV (Yoneda)} must show that the structural role is
captured by the representable functor; in particular, the maps from
$\mathrm{Hom}(-, n)$ to other functors $F$ correspond to elements of
$F(n)$ in a way that respects the everyday operations.
\item \emph{Paper V (HoTT)} must show that the structural role is
captured by the inductive type, with identity types capturing the
intuitive ``same number'' relation native to the naive level.
\item \emph{Paper VI (categorical/structural)} must show that the
structural role is exactly the invariant content under
structure-preserving morphisms.
\end{itemize}

In each case, success of the formal level is measured by recovery of
the naive practice. This is the methodological role of \cref{rem:recovery}.

\subsection{The deflationary alternative}

A reader might worry that we have stacked the deck in favor of
realism about numbers. Could one not, instead, embrace a thoroughly
nominalist view in which numerals are all there is, and arithmetic is
a syntactic discipline of inscription manipulation? The literature
contains serious such proposals (Field, Hellman, in different ways).
We do not engage them in detail here; for our purposes it is enough to
note that any such program must \emph{still} explain why distinct
inscriptions of ``58'' (decimal, Roman, tally) are interchangeable in
arithmetical contexts. That explanation is, in the end, a structural
explanation; the nominalist's parsimony, if successful, is purchased
at the cost of accepting structural facts about inscriptions.

\subsection{Limitations}

We have restricted attention throughout to the natural numbers. The
naive level for the rationals, the reals, the complex numbers, the
$p$-adic numbers, and the algebraic and transcendental numbers raises
analogous but distinct issues. In particular, the relationship between
the naive level and the universal-property level is more delicate for
the real numbers, where the naive level is itself partly definitional
(``what we mean by a real number is a Cauchy-completion / Dedekind cut /
positional infinite expansion''). We defer these to a future companion
piece.

\subsection{The role of computer science}

A subtler limitation: we have spoken of numerals as inscriptions on
paper or stone. In the present era much arithmetic happens in machines,
where numerals are voltage patterns, register configurations, or
floating-point bit fields. The syntax/semantics distinction extends in
the obvious way, with the wrinkle that inexact representations
(floating-point) raise their own issues. We mention this only to flag
that the naive level is not a fixed historical relic; it is a layer of
mathematical practice that travels into new technological substrates
without losing its character.

\section{Conclusion}
\label{sec:conclusion}

We have argued that the naive perspective on number---numerals as syntax,
quantities as semantics, with the two related by counting and
denotation---is at once foundationally inadequate and practically
indispensable. Inadequate, because numbers are not their inscriptions
(\cref{thm:symbol-object}), no inscription is canonical
(\cref{thm:no-canonical-numeral}), and adequate inscriptions are
interconvertible (\cref{thm:naive-benacerraf}). Indispensable, because
acquisition, application, and pedagogy of arithmetic all proceed
through it (\cref{prop:acquisition,prop:application,prop:pedagogical}).

The diagnostic move out of the naive level is the structuralist's
question: \emph{what role does $n$ play?} (\cref{def:structuralist-question}).
The remaining papers in this series formalize this question in five
distinct ways and argue, in the synthesis, that under the univalence
axiom the last four formalizations are equivalent.

\subsection{Future work}

\begin{itemize}
\item Extend the naive analysis to rationals, reals, and complex
numbers, with attention to the role of completion and continuity.
\item Develop the cognitive-science side of the indispensability
arguments more thoroughly: small-number system, approximate number
system, magnitude representations.
\item Pursue the question of inscription on machines (digital,
neuromorphic, quantum), where the naive level is mediated by hardware.
\item Connect the naive level to category-theoretic ``free models'':
on a structuralist reading, the naive level corresponds to working with
the free model on one generator, before any quotient by structural
identifications.
\end{itemize}

\appendix

\section{The Naive Sequence in Detail}
\label{app:naive-sequence}

For completeness we record the first 60 values of the count sequence in
each of our four numeral systems. The point is purely illustrative: the
table makes vivid the shared semantic content of glyphically unrelated
inscriptions.

\begin{center}
\renewcommand{\arraystretch}{1.05}
\begin{tabular}{r|l|l|l}
\textbf{Decimal} & \textbf{Roman} & \textbf{Tally (gate-five)} & \textbf{Sexagesimal} \\ \hline
0  & ---           & (empty)              & 0 \\
1  & I             & |                    & 1 \\
2  & II            & ||                   & 2 \\
3  & III           & |||                  & 3 \\
4  & IV            & ||||                 & 4 \\
5  & V             & ||||/                & 5 \\
\vdots & \vdots    & \vdots               & \vdots \\
10 & X             & two gate-fives       & 10 \\
20 & XX            & four gate-fives      & 20 \\
50 & L             & ten gate-fives       & 50 \\
58 & LVIII         & 11 gate-fives + 3    & 58 \\
59 & LIX           & 11 gate-fives + 4    & 59 \\
60 & LX            & 12 gate-fives        & $1{;}0$ \\
\end{tabular}
\end{center}

The qualitative jump at row 60 in the sexagesimal column---where the
representation acquires a second positional place---illustrates how
positional systems amplify compactness at integer powers of the base.

\section{A Worked Example: Why \texorpdfstring{$58 = \mathsf{LVIII}$}{58 = LVIII} is Non-Trivial}
\label{app:lviii-example}

We work the example through in detail, in part to illustrate the
mechanics of the syntax/semantics distinction, and in part to show how
much of arithmetic the naive theorist already implicitly uses.

\paragraph{Decimal side.} The string ``58'' parses, in the standard
positional algorithm, as $5 \cdot 10 + 8$. To compute this we need the
multiplication $5 \cdot 10 = 50$ and the addition $50 + 8 = 58$. Both
operations are themselves unfolded by the count sequence: we know
$50 + 8 = 58$ because counting from 50 by ones for eight steps gives
51, 52, 53, 54, 55, 56, 57, 58.

\paragraph{Roman side.} The string \textsf{LVIII} parses, by the additive
rule (no leading subtractive in this case), as
$\mathsf{L} + \mathsf{V} + \mathsf{I} + \mathsf{I} + \mathsf{I} =
50 + 5 + 1 + 1 + 1$. We compute $50 + 5 = 55$ (counting from 50 by
ones for five steps) and $55 + 1 + 1 + 1 = 58$ (three more steps).

\paragraph{Equality.} Both sides reduce to 58 (in the count-sequence
sense). The equation $58 = \mathsf{LVIII}$ is therefore the assertion
that two distinct procedures terminate at the same count. The naive
theorist's grasp of this assertion presupposes (a) the count sequence,
(b) the additive principle, (c) confidence that procedures starting
from different inscriptions but tracking the same quantity will agree.
None of these is a primitive of pure inscription manipulation.

\paragraph{Lifting.} The structural lift in subsequent papers will say:
forget the procedures; the equation reflects the fact that both
inscriptions designate the same element of the initial successor
structure / the same Peano term / the same path in $\mathbb{N}$.

\section{On the Word ``Naive''}
\label{app:naive}

We have used ``naive'' throughout in a technical sense; we now record
the choice. Two alternative terms are available: \emph{pre-formal} and
\emph{pre-theoretic}.

``Pre-formal'' is accurate but suggests temporal priority where we
intend logical priority: the naive level is not what comes \emph{before}
formalization in time but what \emph{lies underneath} it, even after
formalization is complete.

``Pre-theoretic'' is misleading because the naive level has plenty of
implicit theory in it: counting, succession, the additive principle. It
is pre-axiomatic, not pre-theoretic.

``Naive'' has the disadvantage of pejorative overtones in everyday
English; we use it in the established mathematical sense (cf.\ ``naive
set theory'', ``naive Bayes'') of \emph{prior to formalization but
serviceable for many purposes}. We hope the reader will tolerate the
term in this technical sense.

\begin{thebibliography}{99}

\bibitem{awodey1996}
S.~Awodey,
\emph{Structure in Mathematics and Logic: A Categorical Perspective},
Philosophia Mathematica \textbf{4} (1996), 209--237.

\bibitem{benacerraf1965}
P.~Benacerraf,
\emph{What Numbers Could Not Be},
The Philosophical Review \textbf{74} (1965), 47--73.

\bibitem{benacerraf1973}
P.~Benacerraf,
\emph{Mathematical Truth},
The Journal of Philosophy \textbf{70} (1973), 661--679.

\bibitem{carey}
S.~Carey,
\emph{The Origin of Concepts},
Oxford University Press, 2009.

\bibitem{dehaene}
S.~Dehaene,
\emph{The Number Sense}, Oxford University Press, 1997.

\bibitem{field1980}
H.~Field,
\emph{Science Without Numbers}, Princeton University Press, 1980.

\bibitem{frege1884}
G.~Frege,
\emph{Die Grundlagen der Arithmetik}, Breslau, 1884.
English translation by J.~L.~Austin, Blackwell, 1953.

\bibitem{frege1892}
G.~Frege,
\emph{\"Uber Sinn und Bedeutung},
Zeitschrift f\"ur Philosophie und philosophische Kritik \textbf{100}
(1892), 25--50.

\bibitem{frege1894}
G.~Frege,
Review of Husserl's \emph{Philosophie der Arithmetik},
Zeitschrift f\"ur Philosophie und philosophische Kritik \textbf{103}
(1894), 313--332.

\bibitem{geier}
M.~Geier (ed.),
\emph{Frege's Lectures on Logic}, Open Court, 2004.

\bibitem{gelman-gallistel}
R.~Gelman and C.~R.~Gallistel,
\emph{The Child's Understanding of Number},
Harvard University Press, 1978.

\bibitem{hale-wright}
B.~Hale and C.~Wright,
\emph{The Reason's Proper Study},
Oxford University Press, 2001.

\bibitem{hellman1989}
G.~Hellman,
\emph{Mathematics Without Numbers},
Oxford University Press, 1989.

\bibitem{husserl1891}
E.~Husserl,
\emph{Philosophie der Arithmetik}, Halle, 1891.

\bibitem{ifrah}
G.~Ifrah,
\emph{The Universal History of Numbers},
Wiley, 2000.

\bibitem{lakoff-nunez}
G.~Lakoff and R.~N\'u\~nez,
\emph{Where Mathematics Comes From}, Basic Books, 2000.

\bibitem{maclane}
S.~Mac~Lane,
\emph{Categories for the Working Mathematician}, 2nd ed., Springer, 1998.

\bibitem{mill1843}
J.~S.~Mill,
\emph{A System of Logic}, London, 1843.

\bibitem{piaget}
J.~Piaget,
\emph{The Child's Conception of Number},
Routledge \& Kegan Paul, 1952.

\bibitem{hottbook}
The Univalent Foundations Program,
\emph{Homotopy Type Theory: Univalent Foundations of Mathematics},
Institute for Advanced Study, 2013. arXiv:1308.0729.

\bibitem{wadler2015}
P.~Wadler,
\emph{Propositions as Types},
Communications of the ACM \textbf{58} (2015), 75--84.

\bibitem{wright1983}
C.~Wright,
\emph{Frege's Conception of Numbers as Objects},
Aberdeen University Press, 1983.

\end{thebibliography}

\end{document}
